\section{Discussion}
\label{sec:hf-discussion}
Due to its practical importance, method-level maintenance research has gained significant traction in recent years, making accurate method-level change history generation increasingly critical. The goal was to assess whether existing tools for generating method-level change histories produce accurate and complete results when evaluated against rigorously constructed oracles, and to identify which tool offers the best balance between accuracy and runtime efficiency. The results show that earlier evaluations based on the original \textit{CodeShovel} and \textit{CodeTracker} oracles were inaccurate, undermining the validity of prior assessments of tool performance (\textbf{RQ1}). To address this, two new oracles—the corrected \textit{CodeShovel} oracle and the new \textit{HistoryFinder} oracle—were systematically constructed, offering a more reliable benchmark that combines automated detection with expert validation.

To further address the limitations observed in both \textit{CodeShovel} and \textit{CodeTracker} tools, a new algorithm and tool, \textit{HistoryFinder}, was also developed to improve both the accuracy and completeness of method change history tracking while maintaining competitive runtime performance. When evaluated against the newly constructed oracles, \textit{HistoryFinder} consistently achieved the highest precision, recall, and F1 scores (\textbf{RQ2}) and demonstrated the most robust performance across all oracles, including those involving challenging repositories with extensive changes to the file containing the target method. Although the \texttt{Git-based} tools delivered faster execution times, this efficiency came at the cost of substantially lower accuracy, making \textit{HistoryFinder} the most reliable option for users prioritizing both accuracy and reasonable runtime (\textbf{RQ3}). The runtime analysis also revealed that repositories with frequent file modifications result in longer processing times for tools that rely on \texttt{JavaParser}.

Both the research community and practitioners can benefit from these contributions. Researchers can leverage these oracles to train and evaluate new algorithms or improve existing method history generation tools, especially in handling complex scenarios—such as frequent file changes, large-scale refactorings that involve method or file moves, and tracking changes in annotations or JavaDocs associated with methods.
A graphical user interface (GUI) is provided that allows both researchers and practitioners to interactively run multiple tools—including \textit{CodeShovel}, \textit{CodeTracker}, \textit{HistoryFinder}, \textit{GitLineRange}, and \textit{GitFuncName}—on their own selected Java methods. Users can independently compare outputs and choose the one best suited to their needs, whether they prioritize accuracy or runtime performance. This tool can also be integrated into a developer’s daily workflow.

In addition to the GUI, this chapter provides support for a command-line interface (CLI) that can be used to track change histories of millions of methods for conducting research on mining software repositories and understanding software evolution at the method level. For example, researchers can explore why some source methods change frequently while others remain stable, and why, when, and how bugs are introduced in source code. Additionally, \textit{HistoryFinder} can potentially be used to track test methods, given the crucial role of test code in software quality and maintainability.
The tool is also shared as a Java library for seamless integration into larger software engineering pipelines. The complete setup and usage instructions are provided with the shared repository.\footnote{\url{https://github.com/SQMLab/HistoryFinder}}

\subsection{Threats to Validity}
\textbf{Internal Validity.}
Although multiple state-of-the-art history generation tools were leveraged to construct the two oracles, the validation of change commits still involved human evaluators. To mitigate this threat, two of the authors—each with over eight years of industry experience—independently conducted the validations. Nonetheless, this threat could be further reduced by involving additional evaluators. Similar to the \textit{CodeShovel} and \textit{CodeTracker} tools, a key limitation of the \textit{HistoryFinder} tool is its reliance on \texttt{JavaParser}, which may fail when processing syntactically incorrect files—often introduced by inadvertent changes during development.

\textbf{External Validity.}
Both the \textit{CodeShovel} and \textit{HistoryFinder} oracles were constructed using only open-source repositories, which may limit the generalizability of the findings. It remains uncertain whether the tools’ performance and accuracy would translate to closed-source or industrial settings, where code characteristics and commit practices may differ. While the \textit{HistoryFinder} tool can be extended to support other programming languages by incorporating language-specific parsers, it is unclear whether the similarity thresholds trained for identifying method matches would remain effective outside of Java.
